% ECONOMICS_PROBLEM: {"id": "D63-294", "title": "EFX existence for chores with binary submodular costs", "jel_code": "D63", "source_id": "OP-0194"}
% !TeX program = xelatex
\documentclass[11pt,a4paper]{article}
\usepackage[margin=23mm]{geometry}
\usepackage{fontspec}
\setmainfont{Latin Modern Roman}
\usepackage{amsmath,amssymb,mathtools}
\usepackage{xcolor,enumitem,booktabs,longtable,array,xurl,hyperref}
\definecolor{muted}{HTML}{53636C}
\hypersetup{colorlinks=true,urlcolor=blue,linkcolor=blue}
\setlength{\parindent}{0pt}
\setlength{\parskip}{5pt}
\newcommand{\R}{\mathbb R}
\newcommand{\E}{\mathbb E}
\newcommand{\N}{\mathbb N}
\newcommand{\ind}{\mathbf 1}
\newcommand{\Var}{\operatorname{Var}}
\newcommand{\Cov}{\operatorname{Cov}}
\newcommand{\supp}{\operatorname{supp}}
\DeclareMathOperator*{\argmin}{arg\,min}
\DeclareMathOperator*{\argmax}{arg\,max}
\newcommand{\entrymeta}[3]{{\sffamily\small\color{muted}\textbf{\texttt{#1}}\quad|\quad #2\par #3\par}}
\newcommand{\Problemref}[1]{\texttt{#1}}
\newcommand{\JELref}[2]{\texttt{#2} (JEL \texttt{#1})}
\begin{document}
\section*{EFX existence for chores with binary submodular costs}
\textbf{Problem ID:} \texttt{D63-294}\quad \textbf{JEL:} \texttt{D63}\par
% BEGIN ECONOMICS STATEMENT
\label{problem:OP-0194}
% GAME_THEORY_PROBLEM: {"local_id": "GT-064", "original_number": 64, "kind": "P", "entry_kind": "statement_only", "published": false, "review_required": true, "category": "game_theory"}
\label{game:GT-064}
{\sffamily\small\color{muted}
\noindent\textbf{ID: GT-064.}\quad\textbf{Category:} Fair division of chores.
\quad\textbf{Kind: P.}\quad\textbf{Status:} Open in the cited source.
\par}
\subsection*{Definitions and mathematical statement}
Let \(N=\{1,\ldots,n\}\), \(n\geq 1\), and let \(C\) be a finite set
of chores. For each \(i\), let
\(c_i:2^C\to\mathbb Z_{\geq 0}\) satisfy \(c_i(\varnothing)=0\),
monotonicity, submodularity, and binary marginal costs:
\[
 \begin{gathered}
 S\subseteq T\ \Longrightarrow\ c_i(S)\leq c_i(T),\\
 c_i(S)+c_i(T)\geq c_i(S\cup T)+c_i(S\cap T)
 \qquad(S,T\subseteq C),\\
 c_i(S\cup\{g\})-c_i(S)\in\{0,1\}
 \qquad(S\subseteq C,\ g\in C\setminus S).
 \end{gathered}
\]
A complete allocation \(A=(A_1,\ldots,A_n)\) is an ordered partition
of \(C\). It is EFX if
\[
 c_i(A_i\setminus\{g\})\leq c_i(A_j)
 \qquad(i,j\in N,\ g\in A_i).
\]

\paragraph{Question.}
Does every instance satisfying these hypotheses admit a complete EFX
allocation?

\subsection*{Short English statement}
Must binary submodular chore costs always permit an EFX allocation?

\subsection*{Variants and scope boundaries}
The removed chore belongs to the potentially envious agent's own bundle.
Every chore is tested, including a chore of zero marginal cost.
Binary marginals do not mean that the value of an entire bundle is binary.
Negative results for other monotone cost classes do not answer the
submodular question.

\subsection*{Source relationship and status}
[S1, Section 4] explicitly singles out binary submodular costs as an
unresolved class.

\subsection*{References}
\begingroup\footnotesize\raggedright
\noindent[S1] Zehan Lin, Shengxin Liu, Biaoshuai Tao, and Shengwei Zhou,
\emph{Non-Existence of EFX Chore Allocations for Monotone Cost Functions
with Binary Marginals}, 2026, arXiv:2608.10572v1.
\url{https://arxiv.org/html/2608.10572v1}.
Locators: Section 2 (cost and EFX definitions);
Section 4, \emph{Conclusion and Open Problem}.
\par\endgroup

\subsection*{Common conventions: Source mathematical conventions}

\textbf{Finite games.} For a finite set $A$, $\Delta(A)$ is its probability
simplex. Players choose independent mixed strategies in Nash equilibrium;
correlation is allowed only when explicitly part of the solution concept.
In a game with payoff functions $u_i$ and strategy profile $x$, an additive
$\varepsilon$ Nash equilibrium satisfies
\[
 u_i(x)\geq u_i(x_i',x_{-i})-\varepsilon
 \quad\text{for every player $i$ and every admissible deviation $x_i'$}.
\]
For numerical approximation, payoffs are normalized as locally specified.
An $\varepsilon$ payoff residual is distinct from distance at most
$\varepsilon$ to the set of exact equilibria. Point convergence, convergence
of empirical play, and convergence in distance to an equilibrium set are
also distinct.

\textbf{Computational representations.} Unless stated otherwise, a complexity
question takes rational data with binary encoding; polynomial time is measured
in total bit length. A fixed-accuracy polynomial algorithm is not necessarily
polynomial in $1/\varepsilon$, and neither is necessarily polynomial in
$\log(1/\varepsilon)$. Succinct games and value, demand, or sampling oracles
must declare their interfaces. Exact real solutions require an explicit
representation; an unspecified list of arbitrary real numbers is not a
finite computational output. A claimed algorithmic barrier is meaningful
only for its specified model and reductions.

\textbf{Dynamic games.} Histories, information, observation, and allowable
randomization are part of the model. A stationary strategy depends only on
the current state; a positional strategy in a deterministic graph game
chooses one action at each controlled vertex. Nash equilibrium at an initial
state and equilibrium after every history are different requirements.
An expectation of a pathwise $\liminf$ payoff, a $\liminf$ of expected averages,
a limiting expected average, and a uniform finite-horizon equilibrium are
different criteria. None is silently substituted for another.

\textbf{Allocation and mechanisms.} A complete allocation assigns every item
exactly once unless otherwise stated. Zero-valued goods, negative values,
capacity constraints, feasibility, lotteries, and copies of goods affect
fairness definitions and are specified locally. Dominant-strategy incentive
compatibility, Bayesian incentive compatibility, universal truthfulness,
and truthfulness in expectation impose different inequalities. Whenever
an objective ratio has zero denominator, its convention is stated or those
instances are excluded.

\textbf{Infinite-dimensional models.} Expectations, integrals, stochastic
equations, and optimization objectives require their local regularity,
integrability, filtrations, and admissible domains. A backward--forward
mean-field system is a boundary-value problem; it is not an initial-value
semiflow without an additional construction. Long-time, large-population,
and vanishing-noise limits require an order of limits and a convergence
topology. If a minimum need not be attained, the objective is its infimum
or supremum and the approximation requirement is stated separately.
% END ECONOMICS STATEMENT
\end{document}
